\documentclass[11pt,a4paper]{article} % Loại bỏ twocolumn, tăng size chữ lên 11pt cho dễ nhìn

% --- Gói lệnh hỗ trợ tiếng Việt và định dạng ---
\usepackage[utf8]{inputenc}
\usepackage[vietnamese]{babel} 
\usepackage[T5]{fontenc}
\usepackage{amsmath, amsfonts, amssymb} 
\usepackage{graphicx} 
\usepackage{cite} 
\usepackage{geometry}
\usepackage{authblk} 
\usepackage{booktabs} 
\usepackage{hyperref} 

% --- Cấu hình lề trang (Dạng 1 cột nên để lề rộng hơn một chút cho chuyên nghiệp) ---
\geometry{top=2.5cm, bottom=2.5cm, left=3cm, right=3cm}

% --- Thông tin bài báo ---
\title{\textbf{Nghiên cứu và Phát triển Hệ thống Giám sát Môi trường Thông minh dựa trên nền tảng IoT và Học sâu}}

\author[1]{Nguyễn Văn A}
\author[2]{Trần Thị B}
\affil[1,2]{Khoa Công nghệ Thông tin, Đại học Bách khoa Hà Nội}
\affil[ ]{Email: \{anv, btt\}@hust.edu.vn}

\date{} 

\begin{document}

\maketitle

\begin{abstract}
Bài báo này trình bày nghiên cứu và thiết kế một hệ thống giám sát môi trường thời gian thực sử dụng công nghệ Internet vạn vật (IoT) kết hợp với các thuật toán học sâu (Deep Learning). Hệ thống cho phép thu thập dữ liệu về nhiệt độ, độ ẩm và nồng độ bụi mịn (PM2.5) từ các trạm cảm biến giá thành thấp, sau đó truyền về máy chủ trung tâm thông qua giao thức MQTT. Chúng tôi đề xuất sử dụng mô hình mạng LSTM (Long Short-Term Memory) để dự báo xu hướng ô nhiễm không khí trong ngắn hạn. Kết quả thực nghiệm cho thấy mô hình đạt độ chính xác cao với sai số trung bình bình phương (MSE) thấp, hứa hẹn khả năng ứng dụng thực tiễn trong các đô thị thông minh.
\end{abstract}

\textbf{Từ khóa:} IoT, Học sâu, LSTM, Giám sát môi trường, MQTT, Đô thị thông minh.

\section{Giới thiệu}
Sự phát triển nhanh chóng của công nghiệp hóa và đô thị hóa đã dẫn đến những thách thức nghiêm trọng về ô nhiễm môi trường. Việc theo dõi chất lượng không khí trở thành nhu cầu cấp thiết để bảo vệ sức khỏe cộng đồng. Các hệ thống quan trắc truyền thống thường có kích thước lớn và chi phí lắp đặt rất cao, gây khó khăn trong việc triển khai diện rộng.

Với sự bùng nổ của cuộc Cách mạng Công nghiệp 4.0, công nghệ IoT cung cấp giải pháp kết nối các thiết bị cảm biến nhỏ gọn với chi phí thấp. Tuy nhiên, dữ liệu thu được thường có độ nhiễu cao. Do đó, việc ứng dụng trí tuệ nhân tạo (AI) để phân tích và xử lý dữ liệu là một hướng đi tất yếu.

\section{Kiến trúc hệ thống}
Hệ thống đề xuất bao gồm ba lớp chính: Lớp thiết bị (Perception Layer), Lớp mạng (Network Layer) và Lớp ứng dụng (Application Layer).

\subsection{Lớp thiết bị}
Chúng tôi sử dụng vi điều khiển ESP32 kết hợp với cảm biến DHT22 (nhiệt độ, độ ẩm) và PMS5003 (bụi mịn). Dữ liệu được lấy mẫu định kỳ mỗi 5 phút một lần. Vi điều khiển sẽ thực hiện tiền xử lý dữ liệu thô trước khi đóng gói gửi đi.

\subsection{Giao thức truyền thông}
Giao thức MQTT (Message Queuing Telemetry Transport) được lựa chọn do đặc tính nhẹ, chiếm ít băng thông và phù hợp cho các thiết bị nhúng có tài nguyên hạn chế. Cấu trúc một gói tin dữ liệu được định nghĩa như sau:
\begin{equation}
    P_{data} = \{ID, Temp, Humid, PM2.5, Timestamp\}
\end{equation}

\section{Mô hình dự báo dựa trên LSTM}
Để dự báo nồng độ bụi mịn trong tương lai, chúng tôi sử dụng mạng LSTM, một biến thể của mạng Nơ-ron tái phát (RNN) có khả năng học các phụ thuộc dài hạn trong dữ liệu chuỗi thời gian.

Cấu trúc của một đơn vị LSTM bao gồm các cổng (gate) để kiểm soát dòng thông tin. Các phương trình toán học cơ bản được mô tả như sau:

\begin{equation}
    f_t = \sigma(W_f \cdot [h_{t-1}, x_t] + b_f)
\end{equation}
\begin{equation}
    i_t = \sigma(W_i \cdot [h_{t-1}, x_t] + b_i)
\end{equation}
\begin{equation}
    \tilde{C}_t = \tanh(W_C \cdot [h_{t-1}, x_t] + b_C)
\end{equation}

Trong đó $\sigma$ là hàm kích hoạt sigmoid, $W$ và $b$ là các trọng số và độ chệch cần được tối ưu hóa. Việc sử dụng LSTM giúp khắc phục hiện tượng triệt tiêu đạo hàm (vanishing gradient) thường gặp trên các mạng RNN truyền thống khi xử lý chuỗi dài.

\section{Kết quả thực nghiệm và Thảo luận}
Hệ thống đã được triển khai thử nghiệm tại khu vực nội thành Hà Nội trong vòng 30 ngày. Tập dữ liệu sau khi làm sạch bao gồm hơn 8,640 mẫu tin hợp lệ.

\subsection{Đánh giá độ chính xác}
Chúng tôi so sánh mô hình LSTM với mô hình hồi quy tuyến tính (Linear Regression) và Rừng ngẫu nhiên (Random Forest). Kết quả được tổng hợp trong Bảng \ref{table:results}.

\begin{table}[h]
\centering
\caption{So sánh hiệu năng giữa các mô hình dự báo}
\label{table:results}
\vspace{0.3cm}
\begin{tabular}{lccc}
\toprule
\textbf{Mô hình} & \textbf{MAE} & \textbf{RMSE} & \textbf{$R^2$ Score} \\
\midrule
Linear Regression & 4.52 & 6.12 & 0.78 \\
Random Forest & 3.15 & 4.28 & 0.85 \\
\textbf{LSTM (Đề xuất)} & \textbf{1.82} & \textbf{2.45} & \textbf{0.94} \\
\bottomrule
\end{tabular}
\end{table}

Kết quả cho thấy mô hình LSTM vượt trội hơn hẳn nhờ khả năng xử lý tính phi tuyến và nắm bắt được tính mùa vụ trong dữ liệu môi trường (ví dụ: nồng độ bụi thường tăng cao vào giờ cao điểm giao thông).

\subsection{Phân tích độ trễ và độ tin cậy}
Độ trễ trung bình từ khi cảm biến thu thập đến khi hiển thị trên Dashboard là khoảng 1.2 giây. Tỷ lệ mất gói tin (Packet Loss Rate) đo được là dưới 0.5\% trong điều kiện mạng 4G ổn định, cho thấy hệ thống hoạt động tin cậy.

\section{Kết luận}
Bài báo đã trình bày một cách tiếp cận toàn diện trong việc xây dựng hệ thống giám sát môi trường thông minh. Việc kết hợp giữa phần cứng IoT linh hoạt và thuật toán học sâu mạnh mẽ đã tạo ra một công cụ hỗ trợ quản lý đô thị hiệu quả. Định hướng tiếp theo của nghiên cứu là tối ưu hóa thuật toán để có thể chạy trực tiếp trên các thiết bị Edge Computing nhằm giảm tải cho máy chủ trung tâm.

\section*{Lời cảm ơn}
Nghiên cứu này được hỗ trợ kinh phí bởi đề tài cấp cơ sở năm 2026 và sự cộng tác của phòng thí nghiệm trọng điểm.

\begin{thebibliography}{99}
\bibitem{ref1} Nguyễn Văn X, "IoT và ứng dụng trong nông nghiệp thông minh", Tạp chí Khoa học và Công nghệ, 2023.
\bibitem{ref2} Hochreiter, S., \& Schmidhuber, J., "Long Short-Term Memory", Neural Computation, 1997.
\bibitem{ref3} Trần Văn Y, "Phân tích dữ liệu chuỗi thời gian", Nhà xuất bản Giáo dục, 2022.
\end{thebibliography}

\end{document}